\chapter{I sistemi e l'innesto}\label{chapter:sistemi}

Questa tesi misura i difetti di recupero di un motore di ricerca piccolo,
Koskidex, dentro un sistema vero, Documentale. Prima di misurare qualsiasi
difetto serve una condizione: che il motore piccolo, messo al posto di
Elasticsearch, trovi gli stessi documenti. Senza quella condizione ogni
miglioramento misurato su Koskidex si potrebbe liquidare con una frase sola:
\emph{stai migliorando il tuo motore, non il nostro sistema}. Questo capitolo
descrive i due sistemi com'erano all'inizio del lavoro, l'innesto dell'uno
nell'altro e le differenze di matching che è stato necessario trovare e
correggere per arrivare agli stessi risultati.

\section{Koskidex: com'era il 23 settembre 2026}\label{sec:koskidex}

Koskidex è un motore di ricerca full-text scritto in Go, nato come progetto
personale per la barra di ricerca di un piccolo negozio online. Si distribuisce
come un binario solo, senza servizi esterni, e l'unica dipendenza del codice è
\texttt{golang.org/x/text}, per la normalizzazione Unicode. Quando è stato scelto
come oggetto di questa tesi contava circa 2.850 righe di codice e 2.060 di test.

\paragraph{Struttura.} Il cuore è un indice invertito in memoria: per ogni
termine la lista delle occorrenze (\emph{posting}), ciascuna con documento,
campo e posizione. Accanto all'indice ci sono i documenti originali, l'elenco
dei termini di ogni documento (per cancellarlo senza scorrere l'indice intero) e
una mappa dai primi due caratteri ai termini che li hanno, usata per cercare i
refusi. La persistenza è un registro delle operazioni (\emph{write-ahead log}) in
JSON, scritto prima di toccare l'indice, più un'istantanea periodica di tutti i
documenti; all'avvio il motore ricarica l'istantanea e riapplica il registro.
Tutto si usa da un'API HTTP: indici, documenti, impostazioni, ricerca con
filtri e faccette.

\paragraph{La ricerca.} Una query si spezza in termini su ogni carattere che non
è una lettera o una cifra, in minuscolo e senza accenti. Ogni termine della
query è \textbf{obbligatorio}: un documento entra fra i risultati solo se li
contiene tutti. Per ogni termine valgono tre forme di corrispondenza:
\begin{itemize}
  \item la forma esatta;
  \item la ricerca per prefisso, per cui \texttt{manut} trova
        \texttt{manutenzione};
  \item i refusi, con la distanza di Damerau~\cite{damerau1964}: uno da quattro
        caratteri in su, due da otto.
\end{itemize}
Il punteggio di un documento è la somma, sui termini, di
\[
  (10 - \text{refusi} + 2 \cdot \text{esatto}) \cdot \text{peso del campo},
\]
dove \emph{esatto} vale 1 se la corrispondenza è esatta. Non c'è frequenza
documentale inversa: un termine che compare in un documento su diecimila pesa
quanto uno che compare in tutti. Un campo \texttt{TF} esisteva nella struttura
delle occorrenze, con il commento \emph{calculated later}, e non era letto da
nessuna parte.

\paragraph{I vettori.} Un documento può avere un vettore, e una ricerca può
portare il vettore della query. In quel caso al punteggio di ogni documento
\emph{già trovato} dalla parte lessicale si aggiunge la similarità del coseno
moltiplicata per 20. È un riordino dei risultati lessicali, non una ricerca
ibrida: un documento che non contiene le parole della query non entra mai.

\paragraph{Il punto di partenza, congelato.} Il comportamento di quel giorno è
fissato in un test, \texttt{TestBaselineRankingIsFrozen}, che confronta gli
ordinamenti di un corpus di prova con quelli registrati e che deve passare,
senza modifiche al test, a ogni commit successivo. Ogni correzione descritta
nei capitoli seguenti sta dietro un'impostazione, con il comportamento di
partenza come valore predefinito. Misurato sulle collezioni pubbliche del
capitolo~\ref{chapter:valutazione}, il punto di partenza fa 0,0246 di nDCG@10 su
SciFact e 0,1659 su NFCorpus.

Letto il codice, quattro scelte sembravano i candidati a difetti di recupero:
il recupero congiuntivo, l'assenza di IDF, il vettore che riordina senza
recuperare e la costante 20 che somma due scale diverse. Sono i difetti dei
capitoli~\ref{chapter:lessicale} e~\ref{chapter:ibrido}. Il primo, come si vedrà,
è stato trovato misurando e non leggendo: la tesi stava per essere scritta su
tre capitoli di ottimizzazione del punteggio di un motore che su SciFact non
restituiva niente in 290 query su 300.

\section{Documentale: come cerca in produzione}\label{sec:documentale}

Documentale è un sistema di gestione documentale sviluppato da Dieffetech: un
backend Laravel, un'interfaccia Next.js, MySQL per i dati ed Elasticsearch per la
ricerca. Chi lo usa archivia atti e documenti in cartelle, con un nome, dei tag,
un riassunto, dei soggetti (gli enti o le persone coinvolte), delle note e dei
metadati liberi.

\paragraph{Due indici.} La ricerca passa da un servizio,
\texttt{ElasticsearchService}, che tiene due indici: uno per i documenti e uno
per le cartelle. L'indice dei documenti ha sei campi di testo, \texttt{name},
\texttt{tags}, \texttt{summary}, \texttt{subjects}, \texttt{notes} e
\texttt{additional\_data}, con l'analizzatore predefinito di Elasticsearch, cioè
il tokenizer standard senza filtri per l'italiano. Il testo integrale dei
documenti non è indicizzato: si cerca sulla scheda.

\paragraph{La query.} La ricerca dei documenti è un \texttt{multi\_match}
di tipo \texttt{best\_fields}~\cite{es-multimatch} con quattro parametri, che il
codice commenta uno per uno:
\begin{itemize}
  \item \texttt{operator: and}, \emph{tutte le parole devono comparire};
  \item \texttt{fuzziness: AUTO}, un refuso da tre caratteri, due da sei;
  \item \texttt{prefix\_length: 1}, la prima lettera non ammette refusi;
  \item i pesi dei campi: \texttt{name} 5, \texttt{tags} 4, \texttt{summary} e
        \texttt{subjects} 3, \texttt{notes} 2, \texttt{additional\_data} 1.
\end{itemize}
Con \texttt{best\_fields} e \texttt{operator: and} tutte le parole della query
devono stare \emph{nello stesso campo}: è un dettaglio che il
capitolo~\ref{chapter:produzione} ritroverà come difetto. La ricerca delle
cartelle è diversa: un \texttt{match} con refusi sul percorso, in disgiunzione,
unito a un \emph{wildcard} \texttt{*termine*}.

\paragraph{Il recupero è congiuntivo anche qui.} Koskidex ed Elasticsearch, in
due basi di codice indipendenti, arrivano alla stessa scelta per la stessa
ragione: in una barra di ricerca, chi scrive due parole si aspetta documenti che
le contengano tutte e due. È la scelta che il capitolo~\ref{chapter:lessicale}
misura come difetto 0 quando le query si allungano.

\paragraph{Il corpus.} Documentale non ha più un archivio di clienti: Dieffetech
lo ha abbandonato nel settembre 2026, e con lui il progetto Elasticsearch in
cloud. Il sistema però gira in locale, con i suoi comandi di importazione e di
indicizzazione, e il contenuto è stato sostituito con un corpus pubblico della
stessa famiglia: 10.018 atti degli albi pretori comunali, 563 del Comune di
Crispiano con il testo integrale e 9.455 di 171 enti del Friuli Venezia Giulia
con i soli metadati. Il corpus è descritto nel
capitolo~\ref{chapter:valutazione}. Qui conta che gli atti entrano in
Documentale dal suo codice di importazione e in Elasticsearch dal suo comando di
indicizzazione, come in produzione: la ricerca misurata è quella dell'app, non
una sua ricostruzione.

\section{L'innesto: un contratto, due motori, gli stessi risultati}\label{sec:innesto}

\paragraph{Il contratto.} L'app non parla con Elasticsearch direttamente ma con
un'interfaccia, \texttt{SearchBackend}, estratta da \texttt{ElasticsearchService}
con una regola sola: le firme sono quelle del servizio in produzione, e se una
firma non torna si corregge l'interfaccia, non il servizio. I metodi sono quelli
che l'app usa davvero: creare i due indici, indicizzare un elemento o un blocco
nel formato \emph{bulk} di Elasticsearch, cancellare, aggiornare, controllare
l'esistenza e cercare (\texttt{fuzzySearch}, che restituisce gli id in ordine di
rilevanza). Accanto a \texttt{ElasticsearchService} c'è \texttt{KoskidexService},
che implementa lo stesso contratto su un client HTTP di Koskidex, e quale dei due
si usa lo decide una variabile d'ambiente, \texttt{SEARCH\_BACKEND}. Nient'altro
dell'app cambia.

\paragraph{Prima dell'innesto, un motore completo.} La regola era di non
innestare un Koskidex che non sapesse fare tutto quello che l'app chiede.
Provandolo contro il contratto sono venuti fuori sette buchi, chiusi uno per uno
in Koskidex con test scritti prima e visti fallire: gli id interi, che venivano
scartati in silenzio; i campi lista (i tag, i soggetti), che venivano salvati ma
non indicizzati; un limite di risultati troppo basso e risposte senza punteggio;
la cancellazione di molti documenti in una richiesta; la ricerca per
sottostringa che serve alle cartelle; e le impostazioni per riprodurre il
matching di Elasticsearch, descritte nella sezione seguente.

\paragraph{Le impostazioni di compatibilità.} \texttt{KoskidexService} crea
l'indice dei documenti con le impostazioni che riproducono la query di
produzione: tutti i termini obbligatori, tutti nello stesso campo, le soglie dei
refusi di \texttt{fuzziness: AUTO}, niente ricerca per prefisso, la prima lettera
esatta, il tokenizer standard e gli stessi pesi dei campi. Il punteggio resta
quello di Koskidex: la compatibilità riguarda \emph{quali} documenti si trovano,
non in che ordine.

\paragraph{Il risultato.} Dal vivo, passando dall'app sui 10.018 atti e cambiando
solo \texttt{SEARCH\_BACKEND}, Koskidex restituisce \textbf{gli stessi insiemi di
Elasticsearch su 24 query su 24} del confronto, e su 300 su 300 delle query
\emph{known-item} del capitolo~\ref{chapter:produzione}, una famiglia di query
che il confronto non conteneva. Gli ordinamenti invece coincidono nei primi dieci
risultati solo su 7 query, quelle con al più due risultati: il punteggio
euristico di Koskidex non è BM25. È esattamente la condizione che serve. Da qui
in poi, ogni differenza misurata fra Koskidex e il suo punto di partenza viene
dalla correzione che si sta misurando, non da un motore che trova altro.

La parità ha un rovescio, e va detto subito: \textbf{riproducendo Elasticsearch,
Koskidex ne eredita i difetti}. Sulle 300 query known-item trova gli stessi
insiemi di Elasticsearch, query per query, e quindi manca gli stessi atti. La parità non è il traguardo ma il punto da
cui misurare le correzioni del capitolo~\ref{chapter:produzione}.

\section{Le differenze di matching trovate per arrivarci}\label{sec:differenze}

La parità non era scontata. Su un primo campione di quattordici documenti con la
forma dei dati di Documentale i due motori trovavano gli stessi insiemi; sul
corpus vero di 10.018 atti coincidevano su 8 query su 24. Il modello di recupero
era lo stesso, congiuntivo, ma il matching differiva. Per trovare in cosa, a
Koskidex è stato dato esattamente il contenuto che Elasticsearch aveva
indicizzato, copiato dall'indice, e si è aggiunta una correzione alla volta,
contando a ogni passo le query con lo stesso insieme e i documenti trovati da un
motore solo (tabella~\ref{tab:divergenza}).

\begin{table}[ht]
  \centering
  \small
  \begin{tabular}{lrrr}
    \toprule
    Passo & identiche & solo ES & solo Koskidex \\
    \midrule
    0. Koskidex com'era & 8 / 24 & 124 & 447 \\
    1. parole nello stesso campo & 10 / 24 & 124 & 121 \\
    2. soglie dei refusi di Elasticsearch & 12 / 24 & 0 & 154 \\
    3. niente ricerca per prefisso & 13 / 24 & 0 & 116 \\
    4. prima lettera esatta & 15 / 24 & 0 & 84 \\
    5. tokenizer standard & \textbf{24 / 24} & \textbf{0} & \textbf{0} \\
    \bottomrule
  \end{tabular}
  \caption{Le differenze fra i due motori sulle 24 query del confronto,
  aggiungendo a Koskidex una correzione alla volta: query con gli stessi risultati,
  e risultati trovati da uno solo dei due motori (ES: Elasticsearch). I passi 2-5 sono misurati
  dopo la correzione del difetto dei bigrammi descritta nel testo.}
  \label{tab:divergenza}
\end{table}

\paragraph{Le cinque cause.}
\begin{enumerate}
  \item \textbf{La congiunzione per campo.} Con \texttt{best\_fields} e
        \texttt{operator: and} tutte le parole devono stare nello stesso campo;
        Koskidex le accettava sparse fra i campi. È la causa più grossa: 326
        dei 447 documenti in più.
  \item \textbf{Le soglie dei refusi.} Elasticsearch ammette un refuso da tre
        caratteri e due da sei, Koskidex da quattro e da otto.
  \item \textbf{La ricerca per prefisso}, che Koskidex fa e Elasticsearch, con
        quella query, no.
  \item \textbf{La prima lettera esatta}, per \texttt{prefix\_length: 1}.
  \item \textbf{Il tokenizer.} Il tokenizer standard di Elasticsearch segue le
        regole di segmentazione di Unicode~\cite{uax29}, che tengono dentro la
        parola un apostrofo fra due lettere e uniscono date e decimali:
        \emph{dell'illuminazione} è un termine solo, e \texttt{14.01.2026}
        anche. Koskidex spezzava su ogni carattere che non è lettera o cifra.
\end{enumerate}
Ciascuna è diventata un'impostazione di Koskidex, spenta per default.

\paragraph{Due ipotesi sbagliate, e un difetto di Koskidex.} Dopo il quarto passo
restavano 9 query diverse, con 42 documenti trovati solo da Elasticsearch, quasi
tutti sulla query \texttt{ordinanza 187}. Le due spiegazioni scritte allora
erano plausibili e sbagliate tutte e due. La prima era la tokenizzazione dei
numeri. La seconda era il limite di espansioni dei refusi di Elasticsearch
(\texttt{max\_expansions}), smentita rilanciando la query di produzione con un
limite di diecimila: stessi insiemi. La causa vera era un difetto di Koskidex.
I candidati per i refusi venivano solo dai termini con almeno un bigramma in
comune con la parola cercata, e \texttt{17} non ne ha nessuno con \texttt{187},
pur distando un refuso. Una modifica rompe fino a due bigrammi, una
trasposizione tre, e sotto una certa lunghezza non resta nessuna garanzia: con
le soglie predefinite capitava anche ad \texttt{atre} contro \texttt{arte}. Il
motore prometteva la distanza di Damerau e non la manteneva. La correzione non è
un'impostazione: per le parole troppo corte perché i bigrammi garantiscano un
candidato, ora si scorre il vocabolario intero, con un filtro sulla lunghezza
prima della distanza.

Il quinto passo è arrivato chiedendo a Elasticsearch, con \texttt{\_explain},
perché non trovava un documento preciso: il nome diceva \emph{dell'illuminazione
pubblica}. È la lezione di metodo di questa sezione, e torna nel resto della
tesi: le ipotesi scritte senza verificarle erano ragionevoli, e l'unica che ha
risolto è venuta da una domanda precisa a un documento preciso.

\paragraph{Un difetto di Documentale, visto per la prima volta.} Il tokenizer
standard è la scelta giusta per le date (\texttt{18} non deve combaciare dentro
\texttt{18.01.2026}) e quella sbagliata per le elisioni: chi cerca
\emph{illuminazione} non trova gli atti che dicono \emph{dell'illuminazione}. Il
capitolo~\ref{chapter:produzione} conta quanti sono.

\section{Prestazioni dell'innesto}\label{sec:prestazioni-innesto}

Con l'innesto fatto, lo stesso confronto dà i costi dei due motori passando
dall'app: indicizzazione dei 10.018 atti da indice vuoto con il comando di
Documentale, e le 24 query tre volte per motore (tabella~\ref{tab:innesto}).

\begin{table}[ht]
  \centering
  \begin{tabular}{lrr}
    \toprule
    & Koskidex & Elasticsearch \\
    \midrule
    insiemi identici sulle 24 query & \multicolumn{2}{c}{24 / 24} \\
    primi dieci identici & \multicolumn{2}{c}{7 / 24} \\
    ricerca, mediana per query & 2,7 ms & 22,0 ms \\
    indicizzazione di 10.018 atti & 1,9 s & 1,7 s \\
    \bottomrule
  \end{tabular}
  \caption{L'innesto misurato dall'app, sugli stessi atti e sulle stesse query.
  Koskidex compilato dal commit \texttt{f583d03}, Elasticsearch 9.1.0 locale.}
  \label{tab:innesto}
\end{table}

I tempi vanno letti come ordini di grandezza, per due ragioni: Elasticsearch
gira nella macchina virtuale di Docker e Koskidex nativo, e tutti e due i tempi
comprendono il viaggio HTTP da PHP. Quello che dicono è che, su un archivio di
diecimila atti, il motore piccolo non costa di più.

\paragraph{La prima misura ha trovato un altro difetto.} La prima
indicizzazione con Koskidex costava 31,8 secondi: 2,6 millisecondi per
documento, costanti. Ogni documento sincronizzava il registro delle operazioni
sul disco, e su macOS una sincronizzazione è un \texttt{F\_FULLFSYNC}, che
attende il supporto fisico. Ora una richiesta scrive tutte le sue righe del
registro con una sincronizzazione sola, e se la scrittura fallisce non entra
niente nell'indice: da 31,8 a 1,9 secondi, con lo stesso formato del registro.

È la seconda volta in questo capitolo, dopo i bigrammi, che la misura trova un
difetto che nessuno stava cercando. Non sarà l'ultima, e non è un caso: un
impianto che misura tutto dall'app, con il codice di produzione, passa dagli
stessi percorsi che percorrerebbe un utente, compresi quelli che nessun test
aveva previsto.
